\section{PHANGS Data and Derived Quantities}
\label{sec:data}

This section describes the observational dataset that anchors the
simulation design. \autoref{sec:phangs_megatable} summarizes the \PHANGS\
megatable survey and the aperture-level join used to construct a single
matched table of gas, stellar, and SFR properties. \autoref{sec:derived}
describes the derivation of the five \PRFM\ design fields from the
\PHANGS\ aperture measurements and stellar/dynamical models.
\autoref{sec:quality_cuts} states the quality cuts applied to define the
reference set $\Dref$ that drives the sampling design in
\CGK{Section~3}. \autoref{tab:fields} summarizes the design and validation
fields with their units and roles.

\subsection{PHANGS Megatable Survey}
\label{sec:phangs_megatable}

We use the publicly released \PHANGS\ megatable products, which compile
spatially resolved measurements of cold gas, ionized gas, stellar mass,
and SFR for a homogeneous sample of $\sim80$ nearby star-forming disk
galaxies \citep{2021ApJS..255...19L,2022AJ....164...43S,2023AJ....166..145L}.
The megatables tile each galaxy in matched apertures sampled on a regular
grid in galactocentric coordinates, and provide for each aperture aperture
photometry of CO(2--1) from \PHANGS-ALMA, \HI\ 21~cm from the \PHANGS-\HI\
companion programs, stellar mass surface density from environmental masks
of Spitzer/IRAC and WISE imaging, and multiple SFR tracers including
H$\alpha$, far-UV (FUV)+WISE-4 (W4) combinations, and an extinction-corrected
H$\alpha$+W4 estimator \citep{2022AJ....164...43S,2023ApJ...945L..19S}.

The megatables are released in two complementary aperture definitions:
hexagonal apertures, which provide a regular tessellation useful for
mapping context and for visual inspection, and two-dimensional
Gaussian-weighted apertures, which provide smoothly tapered measurements
that mitigate edge effects in the aperture sum and are the canonical
aperture for the \PHANGS\ pressure and SFR scaling analyses
\citep{2020ApJ...892..148S,2023ApJ...945L..19S}. \CGK{We adopt the
Gaussian-weighted apertures as canonical throughout this work and use the
hexagonal apertures only for context plots --- please confirm.}

For each galaxy, the gas, stellar, and SFR megatables are released as
separate FITS tables indexed by galaxy and aperture. We perform an inner
join on (galaxy, aperture) across the three tables, retaining only
apertures with valid entries in all joined columns. This yields a single
matched table of approximately $4.6\times10^{4}$ apertures across
$\sim80$ \CGK{galaxies --- please verify the final number after quality
cuts}, which serves as the starting point for the design-field and
validation-field selections described below.

\subsection{Derived Quantities}
\label{sec:derived}

The five environmental parameters that enter \TIGRESS\ as initial
conditions form the design-field vector
\begin{equation}
  \thetaenv \;=\; \left(\Sgas,\, \Sstar,\, \Hstar,\, \Omrot,\, \qshear\right).
  \label{eq:thetaenv}
\end{equation}
We construct each component as follows.

\paragraph{Gas surface density.}
The total neutral gas surface density is the sum of atomic and molecular
contributions,
\begin{equation}
  \Sgas \;=\; \Satom + \Smol,
\end{equation}
where $\Satom$ is derived from the \HI\ 21~cm aperture photometry and
$\Smol$ is derived from the \PHANGS-ALMA CO(2--1) aperture photometry
assuming the metallicity-dependent CO-to-H$_2$ conversion factor adopted
by \citet{2022AJ....164...43S}. Both components include the standard
factor of 1.36 to account for helium. The molecular fraction
$\fmol = \Smol/\Sgas$ is computed from the same components and is used
only as a validation quantity (see below); it is not part of $\thetaenv$.

\paragraph{Stellar surface density.}
The stellar surface density $\Sstar$ is taken from the \PHANGS\ stellar
mass maps based on Spitzer/IRAC 3.6~$\mu$m or WISE-1 3.4~$\mu$m imaging,
calibrated to mass-to-light ratios using independent color information,
following the procedure described in \citet{2021ApJS..255...19L} and
\citet{2025ApJ...989...66V}.

\paragraph{Orbital frequency.}
The orbital angular frequency at each aperture is computed from the
circular velocity of the host galaxy's rotation curve,
\begin{equation}
  \Omrot \;=\; \frac{V_{\rm circ}(R)}{R},
\end{equation}
where $R$ is the galactocentric radius of the aperture center and
$V_{\rm circ}(R)$ is the smooth axisymmetric rotation curve adopted by the
\PHANGS\ dynamical modeling pipeline \citep{2025ApJ...989...66V}.

\paragraph{Stellar scale height.}
The stellar scale height $\Hstar$ is not directly observed and must be
inferred from a stellar dynamical model. We adopt the values tabulated by
\citet{2025ApJ...989...66V}, who fit an axisymmetric stellar mass model
constrained by the rotation curve and the observed stellar surface density
profile, and report the implied vertical scale height under the assumption
that the stellar disk is locally flattened with an axis ratio
characteristic of nearby disk galaxies. \CGK{Please confirm whether
the $\Hstar$ used here is the exponential scale height or the
$\mathrm{sech}^{2}$ scale height in the \citet{2025ApJ...989...66V}
convention; this matters for the mapping to $\zstar$ in
\autoref{sec:tigress_mapping}.}

\paragraph{Shear parameter.}
The shear parameter quantifies the local logarithmic gradient of the
rotation curve and is defined as
\begin{equation}
  \qshear \;\equiv\; -\frac{d\ln \Omrot}{d\ln R}
        \;=\; 1 - \frac{d\ln V_{\rm circ}}{d\ln R}.
  \label{eq:qshear}
\end{equation}
For a flat rotation curve $\qshear = 1$, for solid-body rotation
$\qshear = 0$, and for Keplerian rotation $\qshear = 1.5$. We compute
$\qshear$ aperture-by-aperture from the numerical derivative of the
\PHANGS\ rotation curves at the galactocentric radius of each aperture.
Disk-dominated galaxies are physically bounded by $0 < \qshear \leq 1.5$;
this bound is imposed at the quality-cut stage (\autoref{sec:quality_cuts}).

\subsection{Quality Cuts and Completeness}
\label{sec:quality_cuts}

The reference set $\Dref$ used for sampling is defined by validity cuts
on the five design fields only:
\begin{equation}
  \Dref \;=\; \bigl\{\, \thetaenv \;\big|\;
              \Sgas,\, \Sstar,\, \Hstar,\, \Omrot > 0;\;
              0 < \qshear \leq 1.5 \,\bigr\}.
  \label{eq:dref}
\end{equation}
The positivity requirements on $\Sgas$, $\Sstar$, $\Hstar$, and $\Omrot$
remove apertures with missing or non-finite values in any of the four
quantities. The upper bound $\qshear \leq 1.5$ excludes apertures whose
locally fit rotation-curve derivative is super-Keplerian; such cases occur
near rotation-curve transitions or in noisy outer-disk apertures and are
unphysical for the smooth axisymmetric models that anchor the
\TIGRESS\ initial conditions.

Crucially, we do \emph{not} impose detection cuts on $\fmol$, on the
\HI--CO partition $(\Satom,\Smol)$, or on any of the SFR tracers when
defining $\Dref$. These quantities have their own, generally narrower,
completeness masks: CO is undetected at the lowest gas surface densities,
H$\alpha$ is suppressed in regions of low recent star formation, and W4
saturates in the brightest disks. If validation-field completeness were
imposed on the design selection, the \PHANGS\ prior would be biased toward
high-detection environments---preferentially CO-bright, H$\alpha$-bright,
inner-disk apertures---and the resulting simulation suite would
under-sample the low-pressure, atomic-dominated outer disks that form a
substantial fraction of the total \PHANGS\ disk area. By restricting the
selection cuts to the design fields, the prior reflects the full
disk-area-weighted distribution of galactic environments, and the
validation quantities can be compared against \TIGRESS\ outputs without
introducing a selection bias \CGK{(this argument should be verified against
the actual \PHANGS\ completeness statistics --- please cross-check).}

Validation against \PHANGS\ is performed using only those apertures that
additionally pass the relevant validation-field completeness masks
(detected CO for $\fmol$ and $\Smol$, detected SFR tracer for $\Ssfr$,
and so forth). These masks are applied per-quantity at the validation
stage and are not propagated back into $\Dref$. The full set of design
and validation fields is summarized in \autoref{tab:fields}.

\begin{deluxetable*}{lllp{0.45\textwidth}}
\tablecaption{Design and validation fields used in this work.\label{tab:fields}}
\tablehead{
  \colhead{Field} & \colhead{Symbol} & \colhead{Unit} & \colhead{Role}
}
\startdata
\multicolumn{4}{c}{\textit{Design fields} $\thetaenv$ (selection of $\Dref$)} \\
\hline
Gas surface density        & $\Sgas$  & $\Surf$              & TIGRESS input \\
Stellar surface density    & $\Sstar$ & $\Surf$              & TIGRESS input \\
Stellar scale height       & $\Hstar$ & $\pc$                & TIGRESS input (mapped to $\zstar$) \\
Orbital frequency          & $\Omrot$ & $\kms\,\kpc^{-1}$    & TIGRESS input \\
Shear parameter            & $\qshear$ & dimensionless, $0<\qshear\leq 1.5$ & TIGRESS input \\
\hline
\multicolumn{4}{c}{\textit{Validation fields} (not used for selection of $\Dref$)} \\
\hline
Molecular fraction         & $\fmol$  & dimensionless, $0\leq \fmol \leq 1$ & gas-partition validation \\
Atomic surface density     & $\Satom$ & $\Surf$              & gas-partition validation \\
Molecular surface density  & $\Smol$  & $\Surf$              & gas-partition validation \\
SFR surface density        & $\Ssfr$  & $\sfrunit$           & SFR validation \\
Dynamical equilibrium pressure & $\PDE$ & $\Punit$           & \PRFM\ validation \\
Effective velocity dispersion & $\seff$ & $\kms$             & \PRFM\ validation \\
\enddata
\tablecomments{Design fields are subject only to the validity cuts in
\autoref{eq:dref} and define the reference set $\Dref$ used by the
sampling design. Validation fields carry separate per-quantity
completeness masks and are compared against emergent \TIGRESS\ outputs
in \autoref{sec:validation}.}
\end{deluxetable*}
